% =====================================================================
% BRANCH B DRAFT: Sections 1 and 4
% Working title: When Debiasing Scores Reward Broken Text
% Target: ARR October 2026 cycle (submission Oct 12) -> NAACL/COLING 2027
%
% STATUS: first draft for author review. Citations marked
% [CITATION NEEDED] were NOT programmatically verified and must be
% checked before submission. Verified-real citations are noted inline.
% =====================================================================

\section{Introduction}

Ask a language model how to start a fitness routine and it answers the
question. Mention a disability in the same query and it often answers a
different one: consult your physician, be mindful of your limitations,
consider the risks. The user asked about exercise and was handled as a
patient. This pattern, which we call medicalization bias following
\citep{accesseval2024}, appears across domains that have nothing to do with
medicine, including credit cards, air travel, and film distribution.

Activation steering offers a way to remove such behaviour at inference time,
by editing the model's internal activations with no retraining
\citep{turner2023activation,rimsky2024steering}. Whether an intervention
works is judged by a target metric: how much less often the steered model
produces the vocabulary that marks the bias. We set out to reduce
medicalization bias this way, and the results we report here are what we
found about how that judgment is made.

On the AccessEval disability-bias benchmark \citep{accesseval2024}, the
strongest target-metric score in our comparison belongs to SADI
\citep{wang2025sadi}. At its reported operating point it scores zero, the
metric's value for unbiased text, on 201 of 250 responses. Those responses
are not absent. They average 188 words and read like this: ``Improving
digital security for individuals with mental\texttt{GuidId\&ampHeaderCode}
Behavioral\texttt{HeaderCode} Disorders\ldots'' The medicalization metric is
a log-odds contrast between medical and professional vocabulary, and
corrupted text contains neither, so the contrast evaluates to zero and the
metric records perfect neutrality.

FairSteer \citep{li2025fairsteer} fails the same way at higher strength but
with a different signature: at $\alpha{=}6$ only 68 of its 250 responses are
distinct, and the most frequent one, returned verbatim for 43 different
prompts, is the word ``Here'' repeated for 1{,}048 words. It too scores zero.
Angular Steering \citep{vu2025angular} reaches the highest effect size in our
comparison while injecting corruption inside words, and pushes past neutral,
which withholds medical information from users who asked for it.

These methods sit at the top of a leaderboard. None of them has reduced
medicalization bias.

This is the problem we address. The metric counts an absence, and an
intervention can produce that absence either by removing the bias or by
destroying the answer. Nothing in the standard evaluation separates the two,
and the degenerate outputs are not obviously degenerate: they are fluent in
shape, correctly formatted, and longer than the baseline responses they
replace. A practitioner comparing published numbers would select one of these
methods, deploy it, and ship a model that answers disability-related queries
with corrupted or repeated text. The users who carry that cost are the ones
the intervention was meant to protect.

The same weakness appears a second time, one step earlier in the pipeline.
Recent methods set the \emph{structure} of an intervention, in particular its
rank, from the singular value spectrum of contrastive activation differences
\citep{zou2023representation}. % HiDRA, SAE-SSV, CAS-BiPO: [CITATION NEEDED], all three verified real, BibTeX not yet fetched
A concentrated spectrum is read as a behaviour that one direction can
address, a diffuse one as a behaviour that needs many. We show that this
statistic reports how the contrastive set was built rather than how the
behaviour is represented. Holding the number of pairs fixed at 13 and
varying only the diversity of the prompts moves the top explained-variance
ratio from $22.1\%$ to $76.1\%$, on one benchmark, at one layer.

Both failures have the same shape. A quantity is treated as a measurement of
the bias, and it is measuring something else: output length in the first
case, prompt diversity in the second. We establish this for both quantities
and show what changes once it is corrected. Under a quality-constrained
comparison the ranking of steering methods changes, and a multi-rank
projection reduces medicalization as far as the strongest rank-1 baseline
while keeping the answers intact.

Our contributions:

\begin{itemize}
\item \textbf{An evaluation audit of six steering methods.} We show that
  two of the three highest target-metric scores come from collapsed or
  corrupted generations, and we give the failure modes at token level
  (Section~\ref{sec:audit}).
\item \textbf{A quality-constrained comparison on matched items.} On an
  identical $n{=}250$ AccessEval subset with a shared baseline, the
  multi-rank projection reaches $d{=}0.337$ against CAA's $d{=}0.341$
  while degrading judged quality significantly less (paired difference
  $+0.35$ on a ten-point scale, $95\%$ CI $[0.18, 0.53]$, Wilcoxon
  $p<0.001$).
\item \textbf{Rank matters for the intervention.} Holding the layer, the
  data, and the operator fixed, rank-1 erasure fails ($d{=}{-}0.03$) where
  a rank-40 erasure recovers $d{=}0.205$ (Section~\ref{sec:rank}).
\item \textbf{Measured bias dimensionality is a property of the contrast
  set.} Four controls (prompt homogeneity, cluster versus pair sampling,
  a permutation null, and a cluster bootstrap) show that the spectral
  statistics used to describe ``bias geometry'' track contrast-set
  construction rather than bias type (Section~\ref{sec:geometry}).
\end{itemize}



% =====================================================================
